\documentclass[12pt]{article}

% ── Packages ─────────────────────────────────────────────────────────────────
\usepackage{amsmath, amssymb, amsthm, mathtools}
\usepackage[margin=1in]{geometry}
\usepackage{hyperref}
\usepackage{microtype}
\usepackage{booktabs}
\usepackage{xcolor}
\hypersetup{colorlinks=true, linkcolor=blue!60!black, citecolor=blue!60!black}

% ── Theorem environments ──────────────────────────────────────────────────────
\theoremstyle{definition}
\newtheorem{definition}{Definition}[section]
\newtheorem{remark}[definition]{Remark}

\theoremstyle{plain}
\newtheorem{theorem}{Theorem}[section]
\newtheorem{proposition}[theorem]{Proposition}
\newtheorem{corollary}[theorem]{Corollary}
\newtheorem{lemma}[theorem]{Lemma}

% ── Macros ───────────────────────────────────────────────────────────────────
\newcommand{\Tstar}{T^{*}}
\newcommand{\wstar}{w^{*}}
\newcommand{\muc}{\mu_{c}}
\newcommand{\qfn}{q(w,L,\mu)}
\newcommand{\Pn}{P_{n}(w,\,1-\mu)}
\newcommand{\ceil}[1]{\left\lceil #1 \right\rceil}
\newcommand{\db}{\mathcal{D}}
\newcommand{\R}{\mathbb{R}}

% ─────────────────────────────────────────────────────────────────────────────
\title{Minimum Sequencing Depth Guarantees for\\
       Reference-Based Detection of Divergent Pathogens\\
       in Clinical Metagenomic Sequencing}
\author{Devansh}
\date{\today}
% ─────────────────────────────────────────────────────────────────────────────

\begin{document}
\maketitle

% ── Abstract ─────────────────────────────────────────────────────────────────
\begin{abstract}
We derive the first closed-form lower bound on sequencing depth that
provably guarantees detection of a divergent pathogen in clinical
metagenomic next-generation sequencing (mNGS), jointly as a function
of pathogen relative abundance~$p$, per-base genomic divergence from
the nearest reference~$\mu$, classifier $w$-mer window length~$w$,
read length~$L$, target false-negative probability~$\delta$, and
reference database size~$Q$.
Central to our analysis is the function $\qfn$, giving the
\emph{exact} probability that a single read from a divergent pathogen
produces at least one classifiable $w$-mer match, derived via a
linear recurrence on the longest-run distribution of Bernoulli
trials---a classical combinatorial object appearing here in a
bioinformatics context for the first time.
We further derive the critical divergence~$\muc$: the hard boundary
beyond which no reference-based $w$-mer classifier can guarantee
detection within any fixed sequencing budget, formally separating
detectable from effectively novel pathogens.
Finally, we prove that the optimal classifier window length~$\wstar$
minimizing required depth subject to a false-positive rate
constraint---a sensitivity-specificity tradeoff that grows more acute
as divergence increases---lies at the boundary of the false-positive
feasibility region and is given by an explicit formula.
Our results extend Eidenschink~Brodersen et~al.\ (2020) by introducing
the divergence dimension absent from their model, and provide the
first formal justification for choosing classifier window length from
a detection guarantee rather than empirical benchmarking alone.
\end{abstract}

\tableofcontents
\newpage

% ─────────────────────────────────────────────────────────────────────────────
\section{Model and Definitions}
% ─────────────────────────────────────────────────────────────────────────────

We model a clinical mNGS experiment as follows.
A patient sample contains genomic material from host and microbial
sources.
A pathogen of interest $\mathcal{P}$ is present at \emph{relative
abundance} $p\in(0,1)$: a fraction~$p$ of all sequenced reads
originate from~$\mathcal{P}$.
The pathogen genome has length $N$ base pairs.
The reference database $\db$ contains $Q$ sequences of average
length~$n_{\mathrm{avg}}$ base pairs; the nearest reference
to~$\mathcal{P}$ in~$\db$ diverges from the true genome at per-base
rate~$\mu$ under an \emph{independent} substitution model
(i.e., each base independently disagrees with the reference with
probability~$\mu$).

\begin{definition}[Divergence tier]\label{def:tier}
Let $\mu\in[0,1]$ denote the per-base divergence between a target
pathogen genome and its nearest reference in~$\db$.
\begin{enumerate}
  \item \textbf{Tier~1 (Known):}
        $\mu=0$.  An exact reference exists in~$\db$.
  \item \textbf{Tier~2 (Divergent--detectable):}
        $0<\mu\le\muc$, where $\muc$ is the critical divergence
        defined in Theorem~\ref{thm:mu_c}.
        A reference-based $w$-mer classifier can still guarantee
        $\delta$-detection within a practical sequencing
        budget~$T_{\max}$.
  \item \textbf{Tier~3 (Effectively novel):}
        $\mu>\muc$.
        No reference-based $w$-mer classifier guarantees detection
        within budget~$T_{\max}$.
        Detection requires reference-free methods (de~novo assembly,
        anomaly detection); this regime is out of scope for the
        present work.
\end{enumerate}
\end{definition}

\begin{remark}
$\muc$ is not a fixed biological constant: it depends on~$w$, $L$,
$p$, $\delta$, and~$T_{\max}$.
In particular, increasing read length~$L$ raises~$\muc$ for all
other fixed parameters, providing a formal rationale for the
practical preference for long-read platforms in novel-pathogen
detection.
\end{remark}

\begin{definition}[Informative read]\label{def:informative}
Let all reads have fixed length~$L$ base pairs.
A read~$r$ originating from pathogen~$\mathcal{P}$ with
divergence~$\mu$ is \emph{informative} for a $w$-mer classifier if
it contains at least one substring of length~$w$ that is an exact
match to the corresponding window in the reference sequence.

Under the independent-substitution model, each base of~$r$
independently agrees with the reference with probability $1-\mu$.
Define $P_n = P_n(w,\,1-\mu)$ as the probability that a sequence of
$n$ independent Bernoulli$(1{-}\mu)$ trials contains \emph{no} run
of $w$ consecutive successes.
$P_n$ satisfies the linear recurrence
\begin{equation}\label{eq:recurrence}
  P_n \;=\;
  \sum_{j=0}^{\min(n-1,\,w-1)} (1-\mu)^{j}\,\mu\;P_{n-j-1},
  \qquad n \ge 1,
\end{equation}
with $P_0 = 1$.
The exact informative probability is then
\begin{equation}\label{eq:q_exact}
  \qfn \;=\; 1 - P_L(w,\,1-\mu).
\end{equation}
The Poisson approximation---valid when
$(L-w+1)(1-\mu)^w\ll 1$ and used hereafter for closed-form
expressions---is
\begin{equation}\label{eq:q_poisson}
  \qfn \;\approx\;
  1 - \exp\!\Bigl(-(L-w+1)(1-\mu)^{w}\Bigr).
\end{equation}
\end{definition}

\begin{remark}[Derivation of recurrence~\eqref{eq:recurrence}]
Condition on the position of the first \emph{failure}
(mismatch, probability~$\mu$) in the Bernoulli sequence.
If the first failure occurs at position $j+1$ for
$j\in\{0,\ldots,w-1\}$, the preceding $j$ bases are successes,
no run of length~$w$ has yet been completed, and the problem
reduces to the remaining $n-j-1$ trials.
Summing over all such~$j$ yields~\eqref{eq:recurrence}.
The recurrence has order~$w$ and can be evaluated in $O(wL)$
time, making it computationally tractable for the parameter
ranges of interest ($w\le 40$, $L\le 300$).
\end{remark}

\begin{definition}[Detection event]\label{def:detection}
Given a sequencing run of $T$ reads, let $\mathcal{E}$ denote the
event that at least one informative pathogen read is produced.
The false-negative event is $\overline{\mathcal{E}}$: no informative
pathogen read appears.
We say the classifier \emph{guarantees $\delta$-detection} if
$\Pr(\overline{\mathcal{E}})\le\delta$.
\end{definition}

% ─────────────────────────────────────────────────────────────────────────────
\section{Main Results}
% ─────────────────────────────────────────────────────────────────────────────

\subsection{Minimum Sequencing Depth}

\begin{theorem}[Minimum sequencing depth]\label{thm:depth}
Under the model of Definitions~\ref{def:tier}--\ref{def:detection},
the minimum sequencing depth guaranteeing $\delta$-detection is
\begin{equation}\label{eq:T_star}
  \Tstar(p,\mu,w,L,\delta)
  \;=\;
  \left\lceil
    \frac{\ln(1/\delta)}{-\ln\!\bigl(1-p\cdot\qfn\bigr)}
  \right\rceil
  \;\ge\;
  \left\lceil
    \frac{\ln(1/\delta)}{p\cdot\qfn}
  \right\rceil,
\end{equation}
where the right-hand inequality is the small-$pq$ approximation,
tight to within a factor of~$2$ whenever $p\cdot\qfn\le 0.5$.
\end{theorem}

\begin{proof}
Each read is independently a pathogen read with probability~$p$,
and each pathogen read is independently informative with
probability $q=\qfn$ (Definition~\ref{def:informative}).
Hence each read is an informative pathogen read with
probability~$pq$, independently across reads.
The false-negative event requires all $T$ reads to be
non-informative pathogen reads:
\[
  \Pr(\overline{\mathcal{E}})
  \;=\;
  (1-pq)^{T}.
\]
Setting $\Pr(\overline{\mathcal{E}})\le\delta$ and solving:
\[
  (1-pq)^{T}\le\delta
  \;\iff\;
  T\ln(1-pq)\le\ln\delta
  \;\iff\;
  T\ge\frac{\ln(1/\delta)}{-\ln(1-pq)}.
\]
Taking the ceiling yields the exact minimum integer~$\Tstar$.
The inequality $-\ln(1-x)\ge x$ for $x\in[0,1)$
(with equality only at $x=0$) gives
$-\ln(1-pq)\ge pq$, so
$\ln(1/\delta)/pq\le\Tstar$, establishing the
right-hand approximation.
The factor-of-$2$ claim follows from
$-\ln(1-x)\le 2x$ for $x\le 1/2$.
\end{proof}

\begin{corollary}[Monotonicity]\label{cor:monotone}
$\Tstar$ is strictly increasing in~$\mu$ and strictly increasing
in~$w$, for all $\mu\in(0,1)$, $w\ge 1$, $L>w$.
\end{corollary}

\begin{proof}
$\qfn=1-P_L(w,1-\mu)$ is strictly decreasing in~$\mu$:
larger divergence makes each base more likely to mismatch,
reducing the probability of any run of $w$ consecutive matches.
Similarly, $\qfn$ is strictly decreasing in~$w$: a longer
run is harder to achieve for any fixed $\mu>0$.
Since $\Tstar$ is a strictly decreasing function of $q>0$,
strict monotonicity in both $\mu$ and~$w$ follows.
\end{proof}

\subsection{Critical Divergence}

\begin{theorem}[Critical divergence~$\muc$]\label{thm:mu_c}
Fix $w,L,p,\delta$ and a practical depth ceiling $T_{\max}$.
Define
\[
  \muc(w,L,p,\delta,T_{\max})
  \;=\;
  \sup\bigl\{\mu\in[0,1]\;:\;\Tstar(p,\mu,w,L,\delta)\le T_{\max}\bigr\}.
\]
Under the approximation~\eqref{eq:q_poisson}, $\muc$ is the
unique solution to
\begin{equation}\label{eq:mu_c}
  p\cdot\qfn\Big|_{\mu=\muc}
  \;=\;
  \frac{\ln(1/\delta)}{T_{\max}}.
\end{equation}
For $\mu>\muc$, no reference-based $w$-mer classifier achieves
$\delta$-detection within $T_{\max}$ reads
(Tier~3, Definition~\ref{def:tier}).
\end{theorem}

\begin{proof}
By Theorem~\ref{thm:depth},
$\Tstar\le T_{\max}$ iff
$q(w,L,\mu)\ge\ln(1/\delta)/(pT_{\max})$.
The function $q(w,L,\cdot)$ is continuous and strictly
decreasing in~$\mu$ (Corollary~\ref{cor:monotone}), with
$q(w,L,0)=1-e^{-(L-w+1)}>0$ and $q(w,L,1)=0$.
By the intermediate value theorem there exists a unique
$\muc$ satisfying~\eqref{eq:mu_c}.
For $\mu>\muc$, $q<\ln(1/\delta)/(pT_{\max})$, so
$\Tstar>T_{\max}$, establishing the Tier~3 boundary.
\end{proof}

\begin{remark}
Equation~\eqref{eq:mu_c} can be solved numerically in $O(wL)$
time per evaluation of~$q$ via the recurrence~\eqref{eq:recurrence},
or in closed form using approximation~\eqref{eq:q_poisson}:
\[
  \muc
  \;\approx\;
  1
  -
  \left(
    \frac{-\ln\!\bigl(1-\frac{\ln(1/\delta)}{pT_{\max}}\bigr)}
         {L-w+1}
  \right)^{\!\!1/w}.
\]
\end{remark}

\subsection{Optimal Classifier Window Length}

\begin{theorem}[Optimal $w$-mer length~$\wstar$]\label{thm:optimal_w}
Let $\alpha\in(0,1)$ be a target false-positive rate, where a false
positive arises when a background (non-pathogen) read contains a
$w$-mer matching the pathogen reference.
Modeling database $w$-mer content as uniformly distributed over
$4^w$ possible sequences, the false-positive probability for a
single $w$-mer query is
\begin{equation}\label{eq:fp}
  P_{\mathrm{FP}}(w)
  \;\approx\;
  \frac{Q\cdot n_{\mathrm{avg}}}{4^{w}}.
\end{equation}
The optimal window length minimizing $\Tstar$ subject to
$P_{\mathrm{FP}}(w)\le\alpha$ is
\begin{equation}\label{eq:w_star}
  \wstar
  \;=\;
  \left\lceil
    \log_4\!\left(\frac{Q\cdot n_{\mathrm{avg}}}{\alpha}\right)
  \right\rceil.
\end{equation}
For $\mu>0$, using any $w>\wstar$ increases required sequencing
depth without improving false-positive performance, with the cost
of over-specification diverging as $\mu\to\muc$.
\end{theorem}

\begin{proof}
By Corollary~\ref{cor:monotone}, $\Tstar(w)$ is strictly increasing
in~$w$.
Therefore the minimum of $\Tstar$ over the feasible set
$\mathcal{W}=\{w\in\mathbb{Z}^+:P_{\mathrm{FP}}(w)\le\alpha\}$
is attained at $w^{\dagger}=\min\mathcal{W}$.
The constraint $P_{\mathrm{FP}}(w)\le\alpha$
with~\eqref{eq:fp} becomes
\[
  4^{w}\ge\frac{Q\cdot n_{\mathrm{avg}}}{\alpha}
  \;\iff\;
  w\ge\log_4\!\left(\frac{Q\cdot n_{\mathrm{avg}}}{\alpha}\right),
\]
so $w^{\dagger}=\lceil\log_4(Q\cdot n_{\mathrm{avg}}/\alpha)\rceil=\wstar$.

The cost of over-specification: for $w=\wstar+\Delta$ with $\Delta\ge 1$,
\[
  \frac{\Tstar(\wstar+\Delta)}{\Tstar(\wstar)}
  \;=\;
  \frac{q(\wstar,L,\mu)}{q(\wstar+\Delta,L,\mu)}
  \;\xrightarrow{\mu\to\muc}\;
  \infty,
\]
since $q(\wstar,L,\muc)=\ln(1/\delta)/(pT_{\max})>0$ while
$q(\wstar+\Delta,L,\muc)\to 0$ strictly faster
(by strict monotonicity in~$w$).
\end{proof}

\begin{remark}[Tradeoff sharpening under divergence]
At $\mu=0$, $q(w,L,0)=1-e^{-(L-w+1)}\approx 1$ for all
$w\ll L$, so $\Tstar$ is nearly independent of~$w$ and any
$w\ge\wstar$ is acceptable.
For $\mu>0$, $\Tstar$ grows strictly in~$w$ beyond~$\wstar$,
and the penalty is largest precisely when the pathogen is most
divergent---i.e., when correct window choice matters most.
This is the central design implication: as genomic divergence
increases, the penalty for over-specifying~$w$ grows, making
the principled choice of $w=\wstar$ progressively more
consequential.
\end{remark}

% ─────────────────────────────────────────────────────────────────────────────
\section{Numerical Calibration}
% ─────────────────────────────────────────────────────────────────────────────

Table~\ref{tab:calibration} shows representative values of
$\Tstar$ and $\wstar$ across biologically relevant parameter
ranges, computed using the Poisson approximation~\eqref{eq:q_poisson}.

\begin{table}[h!]
\centering
\renewcommand{\arraystretch}{1.3}
\begin{tabular}{@{}ccccc|cc@{}}
\toprule
$p$ & $\mu$ & $w$ & $L$ & $\delta$ & $\Tstar$ (approx.) & Notes\\
\midrule
$10^{-4}$ & $0.00$ & 24 & 150 & 0.01 & $\sim\!46{,}000$ & Tier~1 \\
$10^{-4}$ & $0.05$ & 24 & 150 & 0.01 & $\sim\!60{,}000$ & Tier~2, mild divergence \\
$10^{-4}$ & $0.15$ & 24 & 150 & 0.01 & $\sim\!200{,}000$ & Tier~2, high divergence \\
$10^{-6}$ & $0.02$ & 24 & 150 & 0.01 & $\sim\!4.6\times 10^6$ & mNGS for rare pathogen \\
$10^{-6}$ & $0.02$ & 30 & 150 & 0.01 & $\sim\!11\times 10^6$ & over-specified $w$ \\
\bottomrule
\end{tabular}
\caption{Representative depth bounds.
$Q=5{,}000$, $n_{\mathrm{avg}}=15{,}000$~bp,
$\alpha=10^{-6}$ giving $\wstar=24$.
The last row illustrates the cost of choosing $w=30>\wstar$.}
\label{tab:calibration}
\end{table}

For $Q=5{,}000$, $n_{\mathrm{avg}}=15{,}000$~bp,
$\alpha=10^{-6}$:
\[
  \wstar
  = \left\lceil\log_4\!\left(\frac{5000\times 15000}{10^{-6}}\right)\right\rceil
  = \left\lceil\log_4(7.5\times 10^{13})\right\rceil
  = \lceil 23.1 \rceil
  = 24.
\]
At $L=150$, $w=24$, $\mu=0.02$:
$\lambda=(150-24+1)\times 0.98^{24}=127\times 0.617\approx 78$,
giving $q\approx 1-e^{-78}\approx 1.00$,
confirming Tier~2 detection is not depth-limited by divergence at
typical sequencing depths for this parameter regime.

% ─────────────────────────────────────────────────────────────────────────────
\section{Limitations and Scope}
% ─────────────────────────────────────────────────────────────────────────────

The results above are derived under an idealized Level~1 model.
Several real-world complications require extension in future work:

\paragraph{Non-uniform coverage.}
The model assumes each read position is drawn uniformly over the
pathogen genome; GC bias and amplification artifacts violate this.

\paragraph{Sequencing errors.}
Error rate~$e$ can be incorporated by replacing~$\mu$ with
$\mu_{\mathrm{eff}}=\mu+e-\mu e$ throughout, increasing $\Tstar$
by at most a factor of $(1-e)^{-w}$ relative to the error-free case.

\paragraph{Host depletion.}
In practice, host reads are the majority of a clinical sample;
$p$ here refers to the post-depletion abundance.
The bound applies to the effective post-depletion~$p$.

\paragraph{Truly novel pathogens ($\mu>\muc$).}
Theorem~\ref{thm:mu_c} formally marks this boundary.
Reference-free detection methods---and their own sample-complexity
theory---are required beyond it.

\end{document}
